\documentclass[11pt]{article}
\usepackage[margin=1in]{geometry}
\usepackage{amsmath,amssymb,amsthm,mathtools}
\usepackage[ruled,vlined,linesnumbered]{algorithm2e}
\usepackage{hyperref}
\usepackage{xcolor}
\usepackage{enumitem}
\usepackage{booktabs}

\hypersetup{colorlinks=true,linkcolor=blue!50!black,citecolor=blue!50!black,urlcolor=blue!50!black}

% --- theorem-like environments, mirroring the outline's own numbering style
\theoremstyle{plain}
\newtheorem{thm}{Theorem}[section]
\newtheorem{cor}[thm]{Corollary}
\newtheorem{prop}[thm]{Proposition}
\newtheorem{lem}[thm]{Lemma}
\theoremstyle{definition}
\newtheorem{defn}[thm]{Definition}

% --- a distinct, clearly-marked verbatim-quote block, so quoted material is
%     never typographically confusable with this document's own claims.
\newenvironment{paperquote}[1]{%
  \def\paperquotecite{#1}%
  \begin{quotation}\noindent\itshape\small
}{%
  \normalfont\par\hfill---\,\textsc{Outline}, \paperquotecite\medskip
  \end{quotation}
}
\newcommand{\pq}[2]{\begin{paperquote}{#1}#2\end{paperquote}}

% --- a marker for content that is MY supporting claim, not a paper quote
\newenvironment{mynote}{%
  \begin{quotation}\noindent\small
}{%
  \end{quotation}
}

\title{Computational Verification of BB Code Automorphisms:\\
\large A Step-by-Step Correspondence to\\
\textit{Automorphisms of Bivariate Bicycle Codes via Refined CRT\\
Components and Local Transfer Data} (outline, AD, HS, JB, IC)}
\author{Companion document to \texttt{bb\_automorphisms} toolkit}
\date{}

\begin{document}
\maketitle

\begin{abstract}
This document is a line-by-line correspondence between the computational
forward-detection and inverse-design algorithms implemented in the
\texttt{bb\_automorphisms} toolkit and the exact theorem, definition, and
corollary statements of the outline paper that motivate them. Every
paper-derived claim is set in a quotation block, attributed by section or
theorem number, and followed immediately by the algorithm step that
computationally realizes it. Passages set as \emph{notes} (unboxed) are the
author's own supporting claims used to justify a computational shortcut;
they are not verbatim statements from the outline and are flagged as such.
\end{abstract}

\tableofcontents

%==============================================================================
\section{Algebraic preliminaries quoted from the outline}
%==============================================================================

Every algorithm below operates on one of a small number of algebraic
objects, all defined directly in the outline. We collect the exact
definitions here so that later quotations can refer back to them without
repetition.

\subsection{The ambient ring $R_{\ell,m}$}

\pq{\S2.2, eqs.\ (1)--(2)}{%
``To understand BB codes we must first understand their ambient ring. Where
GB codes live in $R_\ell^2$, BB codes instead live in $R_{\ell,m}^2$ where
$R_{\ell,m} := R_\ell \otimes R_m$, and
\[ R_\ell := \mathbb{F}_2[x]/\langle x^\ell - 1\rangle \]
(and similarly for $R_m$, w.r.t the variable $y$).''
Expanding: ``$R_{\ell,m} = \mathbb{F}_2[x,y]/\langle x^\ell-1,y^m-1\rangle
\cong \mathbb{F}_2[x]/\langle x^\ell-1\rangle \otimes
\mathbb{F}_2[y]/\langle y^m-1\rangle = R_\ell \otimes R_m$.''
}

\pq{\S2.2, eq.\ (3)}{%
``Alternatively, we can view $R_{\ell,m}$ through its group algebra:
$R_{\ell,m} \cong \mathbb{F}_2[C_\ell \times C_m]$, where ideals are
bivariate abelian codes.''
}

\begin{mynote}
\textbf{Computational representation.} A polynomial in $R_{\ell,m}$ is
represented as a dense $\ell\times m$ array over $\mathbb{F}_2$ (class
\texttt{bb\_ring.Ring}), indexed by the monomial basis
$\{x^iy^j : 0\le i<\ell,\ 0\le j<m\}$ exactly as in eq.\ (2)--(3) above; ring
multiplication is circular (group-algebra) convolution.
\end{mynote}

\subsection{The Fisher--Sehgal trichotomy}

\pq{Thm.\ 2.1}{%
``Theorem 2.1 (Fisher-Sehgal [7]). A group algebra $\mathbb{F}_p^k[G]$ is a
principal ideal group algebra (PIGA) if and only if the Sylow $p$-subgroup
of $G$ is cyclic.''
}

\pq{\S2.2}{%
``With $G = C_\ell \times C_m$, the Sylow 2-subgroup is
$C_{2^{s_\ell}} \times C_{2^{s_m}}$ and we arrive at the following
trichotomy: [\ldots] in the case that both $\ell,m$ even, provably at least
one ideal is not principal.''
\par\smallskip
\begin{tabular}{@{}lll@{}}
\toprule
Regime & Sylow 2-subgroup & PIGA? \\
\midrule
$\ell,m$ odd   & trivial & yes (and semisimple) \\
$\ell,m$ mixed & $C_{2^s}$ cyclic & yes \\
$\ell,m$ even  & $C_{2^{s_\ell}}\times C_{2^{s_m}}$ non-cyclic & no \\
\bottomrule
\end{tabular}
}

\subsection{The refined CRT decomposition and its ``H0'' property}

\pq{Cor.\ 2.5, ``(H0)''}{%
``Corollary 2.5 ((H0)). $R_{\ell,m} \cong \prod_{c\in C} J_c$ with every
$J_c$ finite local, with no dependance on $\ell,m$ beyond being finite
non-negative integers.''
}

This single fact --- that $R_{\ell,m}$, whether or not it is a principal
ideal ring (Thm.\ 2.1), is \emph{always} a finite product of \emph{local}
rings --- is what licenses every Nakayama-type argument used later, in both
the $n=1$ (Cor.\ 3.6) and $n$-generated (Thm.\ 3.2) cases; see \S\ref{sec:shared}.

\subsection{Bar involution, circulant lift, and code parameters}

\pq{\S2.3.1}{%
``$\overleftarrow{f}(x,y) := f(x^{-1},y^{-1})$'' \quad and \quad
``Define the circulant matrix lifting operators:
$\rho : x^iy^j \mapsto \mathrm{circ}_\ell(x^i)\otimes\mathrm{circ}_m(y^j)$.
Setting $A=\rho(f_1)$, $B=\rho(f_2)$, we obtain the $H_X,H_Z$ parity check
matrices: $H_X=[A|B]$, $H_Z=[B^T|A^T]$.''
}

\pq{\S2.3.1}{%
``Such a code is a $[[2\ell m, k]]$ code, where
$k = 2\dim(R_{\ell,m}/\langle f_1,f_2\rangle) =
2\sum_{c\in C}\dim(J_c/\langle f_{1,c},f_{2,c}\rangle)$.''
}

\begin{mynote}
\texttt{bb\_code.BBCode} builds $H_X,H_Z$ exactly via this $\rho$
(\texttt{\_circulant\_lift}), and computes $k$ as $n-\mathrm{rank}(H_X)-
\mathrm{rank}(H_Z)$ over $\mathbb{F}_2$ --- the direct GF(2)-linear-algebra
form of the dimension count above, without evaluating the right-hand
component sum $\sum_c\dim(J_c/\langle\cdot\rangle)$ explicitly.
\end{mynote}

\subsection{The 2-generated stabilizer module}

\pq{\S2.3.1, eqs.\ (29)--(30)}{%
``$\mathrm{StabSpace}(C) = \{(a\cdot f_1,\, a\cdot f_2,\, b\cdot
\overleftarrow{f_2},\, b\cdot\overleftarrow{f_1}) \mid a,b\in R_{\ell,m}\}$
[\ldots] This yields an $R_{\ell,m}\oplus R_{\ell,m}$ module, generated by
$g_X=(f_1,f_2,0,0)$ and $g_Z=(0,0,\overleftarrow{f_2},\overleftarrow{f_1})$.''
}

\subsection{The ring-automorphism catalog}

\pq{\S2.2}{%
``Every group automorphism $\gamma\in\Gamma:=\mathrm{Aut}(G)$ extends to an
$\mathbb{F}_2$-algebra automorphism of $R_{\ell,m}$; these are the ring
automorphisms realizable as permutations of the $\ell\times m$ qubit grid.
The catalog:
\begin{enumerate}[label=\arabic*.,itemsep=2pt]
\item Multipliers $\psi_{(j_x,j_y)}$: $x\mapsto x^{j_x}$, $y\mapsto
  y^{j_y}$, $(j_x,j_y)\in(\mathbb{Z}/\ell)^\times\times(\mathbb{Z}/m)^\times$.
\item Partial folds $\theta_x: x\mapsto x^{-1}, y\mapsto y$ and $\theta_y$:
  involutions folding one axis, with no univariate analogue.
\item The full fold $\iota=\theta_x\theta_y$, i.e.\ the antipode, the
  multiplier $(-1,-1)$.
\item Shears: $\mathrm{Hom}(\mathbb{Z}_\ell,\mathbb{Z}_m)\cong\mathbb{Z}_d$,
  $d=\gcd(\ell,m)$; for $d>1$, $x\mapsto xy^{sm/d}$, $y\mapsto y$ (and
  transposes) [\ldots]
\item The transpose $\tau: x\leftrightarrow y$ when $\ell=m$.
\end{enumerate}''
}

\pq{\S2.2}{%
``Shifts (multiplication by $x^iy^j$) are deliberately excluded: they are
$R_{\ell,m}$-module automorphisms, not ring automorphisms\ldots''
}

\begin{mynote}
\texttt{bb\_ring.py} implements all five catalog families as explicit maps
$(i,j)\mapsto(i',j')$ on exponents (functions \texttt{multiplier},
\texttt{theta\_x}, \texttt{theta\_y}, \texttt{full\_fold},
\texttt{shear\_family} \emph{and} \texttt{shear\_family\_transpose} for the
parenthetical ``(and transposes)'' above, \texttt{transpose\_auto}).
Item 4's transpose clause is easy to under-implement --- an earlier version
of this toolkit omitted \texttt{shear\_family\_transpose} entirely; see
\S\ref{sec:shared}.
\end{mynote}

%==============================================================================
\section{The computational bridge: row space $=$ cyclic submodule}
\label{sec:bridge}
%==============================================================================

\begin{mynote}
This section states, precisely, the one non-paper-quoted fact that the
entire computational approach rests on. It is not verbatim in the outline,
but it is an immediate consequence of the group-algebra structure the
outline itself uses throughout \S2.2--\S3.
\end{mynote}

\begin{lem}[Rowspace realizes the cyclic submodule]
\label{lem:rowspace}
For $R=\mathbb{F}_2[G]$ a group algebra and $g\in N$ (a free $R$-module),
\[ R\cdot g \;=\; \mathrm{span}_{\mathbb{F}_2}\{s\cdot g : s\in G\}. \]
\end{lem}
\begin{proof}
Every $r\in R$ is an $\mathbb{F}_2$-linear combination of group elements
$s\in G$ (this is the definition of the group algebra), and the action of
$R$ on $N$ is $\mathbb{F}_2$-linear, so $r\cdot g$ is the corresponding
linear combination of the $s\cdot g$. Conversely each $s\cdot g$ is $r\cdot
g$ for $r=s\in R$.
\end{proof}

Applied to $g=(f_1,f_2)\in N=R_{\ell,m}^2$: the matrix $H_X=[A|B]$ has one
row per group element $s=x^iy^j\in G=\mathbb{Z}_\ell\times\mathbb{Z}_m$
(the row indexed by $(i,j)$ is $s\cdot(f_1,f_2)$, exactly the circulant
shift construction of \S2.3.1 quoted above), so by Lemma~\ref{lem:rowspace},
\[ \mathrm{rowspace}(H_X) \;=\; R_{\ell,m}\cdot(f_1,f_2). \]
That is, $H_X$'s row space \emph{is} the cyclic submodule $M_X$ that Thm.\
3.1 and Cor.\ 3.6 (quoted in \S\ref{sec:forward} below) are statements
about --- not an encoding of it, the object itself, as an
$\mathbb{F}_2$-vector space. This is what lets every module-equality
question below be answered by ordinary $\mathrm{GF}(2)$ Gaussian
elimination.

%==============================================================================
\section{Forward problem: automorphism detection}
\label{sec:forward}
%==============================================================================

\subsection{Algorithm F1 --- ring-level detection (Corollary 3.6)}

\pq{\S3.1}{%
``(H1) Semilinearity: There exists a ring automorphism $\psi\in
\mathrm{Aut}(R_{\ell,m})$ such that, for all $r\in R_{\ell,m}, v\in N$,
$\sigma(r\cdot v)=\psi(r)\cdot\sigma(v)$.
(H2) Cyclic source: $M=R_{\ell,m}\cdot g$ for a single generator element
$g\in N$\ldots''
}

\pq{Thm.\ 3.1}{%
``Assume (H1), (H2). As $R_{\ell,m}$ is local at every component, we have
that
\[ \sigma(M)=M \iff \sigma(g)=w\cdot g \text{ for } w\in R_{\ell,m}^\times.
\]
Proof (Sketch). (H0) means that $\sigma(M)=R\cdot\sigma(g)$ is again
cyclic, reducing to a question about when two cyclic modules with given
generators coincide. At each component: $M_c=J_c\subseteq N_c$, with $g_c$
the image of $g$. Apply Nakayama at each $c$, which we can do as $J_c$ is
local even if $R_{\ell,m}$ is not\ldots Hence $r_c\in J_c^\times$ at every
$c$, which turns a unit at every component into a global unit.''
}

\pq{Cor.\ 3.6}{%
``Given a BB code $C$ with generator $(f_1,f_2)$ over $R_{\ell,m}$, $\phi$
is a code automorphism of the code if $\phi((f_1,f_2))=u(f_1,f_2)$, where
$u$ is a unit in $R_{\ell,m}^\times$.''
}

\begin{mynote}
Algorithm F1 decides the right-hand side of Thm.\ 3.1/Cor.\ 3.6's
``$\iff$'' \emph{without} constructing the component decomposition $J_c$ or
the unit $w$ itself; it decides the left-hand side directly, via
Lemma~\ref{lem:rowspace}, and the theorem's own equivalence licenses
treating that as settling the right-hand side too.
\end{mynote}

\begin{algorithm}[H]
\caption{Ring-level automorphism test (realizes Thm.\ 3.1 / Cor.\ 3.6)}
\SetKwInOut{Input}{Input}\SetKwInOut{Output}{Output}
\Input{code $(\ell,m,A,B)$; candidate $\psi\in$ catalog; $\mathrm{swap}\in\{0,1\}$}
\Output{\textsc{true} iff $\psi$ (composed with the slot-swap) is a code automorphism}
$A' \gets \psi(A)$ \tcp*{apply $\psi$'s exponent-lattice map to every monomial of $A$}
$B' \gets \psi(B)$\;
\If{$\mathrm{swap}$}{swap $A'\leftrightarrow B'$ \tcp*{tests $\phi((f_1,f_2))\propto(f_2,f_1)$ instead of $(f_1,f_2)$}}
$H_X' \gets [\rho(A')\,|\,\rho(B')]$ \tcp*{same construction as $H_X$, from \S2.3.1}
$r_1,r_2,r_3 \gets \mathrm{rank}_{\mathbb F_2}(H_X),\ \mathrm{rank}_{\mathbb F_2}(H_X'),\ \mathrm{rank}_{\mathbb F_2}\!\begin{pmatrix}H_X\\H_X'\end{pmatrix}$\;
\Return{$r_1=r_2=r_3$} \tcp*{$\iff \mathrm{rowspace}(H_X)=\mathrm{rowspace}(H_X') \iff \psi(f_1,f_2)=w(f_1,f_2),\ w\in R^\times$ \textbf{[Thm.\ 3.1]}}
\end{algorithm}

\subsection{Algorithm F2 --- full stabilizer-level detection (Theorem 3.2)}

\pq{Thm.\ 3.2}{%
``Assume (H1) holds. $R_{\ell,m}$ is still a per-component local ring. Let
$M$ be generated by an $n$-tuple $g=(g_1,\ldots,g_n)$ of elements of $N$.
Then
\[ \sigma(M)=M \iff \exists\, g_{\mathrm{block}}\in M_n(R_{\ell,m}),\
\sigma(g)=g_{\mathrm{block}}\,g,\ \text{and } g_{\mathrm{block}}\text{
invertible at every } c. \]''
}

\pq{\S3.1.1}{%
``$g_{\mathrm{ring}}$ elements act as coordinate permutations on a single
$R_{\ell,m}$ element\ldots $g_{\mathrm{ring}}$ elements never act on two
$R_{\ell,m}$ elements at once. [\ldots] $g_{\mathrm{block}}$ elements on the
other hand govern interactions between $R_{\ell,m}^b$ module entries and
encode transversal, fold transversal, and other low-depth unitary
transformations\ldots''
}

The outline restricts its own worked examples to \emph{scalar} (0/1)
$g_{\mathrm{block}}$ entries, i.e.\ $g_{\mathrm{block}}\in M_4(\mathbb{F}_2)$
rather than the fully general $M_4(R_{\ell,m})$:

\pq{\S3.1.2--\S3.1.4, eqs.\ (34)--(37)}{%
Swap-fold: $\begin{psmallmatrix}0&I&0&0\\I&0&0&0\\0&0&0&I\\0&0&I&0\end{psmallmatrix}$
(``corresponds to swapping on the $(i,i+\ell m)$ fold'').
CX-fold: $\begin{psmallmatrix}I&I&0&0\\0&I&0&0\\0&0&I&0\\0&0&I&I\end{psmallmatrix}$
(``corresponds to the action of a CNOT controlled on $i$, targeting
$i+\ell m$'').
Hadamard-type (two variants): $\begin{psmallmatrix}0&0&I&0\\0&0&0&I\\I&0&0&0\\0&I&0&0\end{psmallmatrix}$
and $\begin{psmallmatrix}0&0&0&I\\0&0&I&0\\0&I&0&0\\I&0&0&0\end{psmallmatrix}$.
CZ-fold: $\begin{psmallmatrix}I&0&0&0\\0&I&0&0\\I&0&I&0\\0&I&0&I\end{psmallmatrix}$
(``CZ's along the $x^{-1}y^{-1}$ fold'').
}

\begin{algorithm}[H]
\caption{Full stabilizer-module symmetry test (realizes Thm.\ 3.2, restricted to $\psi_L=\psi_R$ and scalar $g_{\mathrm{block}}$)}
\SetKwInOut{Input}{Input}\SetKwInOut{Output}{Output}
\Input{code; candidate $\psi\in$ catalog; $g_{\mathrm{block}}\in\{$eqs.\ (34)--(37), identity$\}$}
\Output{\textsc{true} iff $g_{\mathrm{block}}\circ(\psi\oplus\psi\oplus\psi\oplus\psi)$ preserves $M$}
$S \gets \mathrm{blockdiag}(H_X,H_Z)$ \tcp*{$2\ell m\times 4\ell m$; row space $=M=R{\cdot}g_X\oplus R{\cdot}g_Z$ [eq.\ (29)--(30)]}
\For{$k\gets 1$ \KwTo $4$}{permute column-block $k$ of $S$ by $\psi$'s grid map $\to S_1$ \tcp*{$\Phi=\psi\oplus\psi\oplus\psi\oplus\psi$, the $\psi_L=\psi_R$ case of $\S3.1$}}
\ForEach{output block $j\in\{1,\ldots,4\}$}{block $j$ of $S_2 \gets \bigoplus_{i:\,g_{\mathrm{block}}[j,i]=1}$ (block $i$ of $S_1$)\;}
$r_1,r_2,r_3 \gets \mathrm{rank}(S),\ \mathrm{rank}(S_2),\ \mathrm{rank}\!\begin{pmatrix}S\\S_2\end{pmatrix}$\;
\Return{$r_1=r_2=r_3$} \tcp*{$\iff \exists\,g_{\mathrm{block}}$ as in Thm.\ 3.2 with $\sigma(g)=g_{\mathrm{block}}g$}
\end{algorithm}

\subsection{Algorithm F3 --- catalog construction and scan}

\pq{\S2.2 (quoted in full above)}{See the five-item catalog list in \S1.6.}

\begin{algorithm}[H]
\caption{Catalog construction (BFS closure) and full scan}
\SetKwInOut{Input}{Input}\SetKwInOut{Output}{Output}
\Input{$(\ell,m)$; code}
\Output{all detected symmetries at both levels}
$\mathcal{G} \gets \{$all multipliers$\}\cup\{\theta_x,\theta_y,\iota\}\cup\{$shear, shear-transpose families$\}\cup\{\tau$ if $\ell=m\}$ \tcp*{catalog item 1--5, \S2.2}
$\Gamma \gets$ BFS closure of $\langle\mathcal{G}\rangle$ in $\mathrm{Sym}(\mathbb{Z}_\ell\times\mathbb{Z}_m)$, deduplicated by resulting permutation\;
\ForEach{$\psi\in\Gamma$, $\mathrm{swap}\in\{0,1\}$}{run Algorithm F1$(\psi,\mathrm{swap})$\;}
\ForEach{$\psi\in\Gamma$, $g_{\mathrm{block}}\in\{$eqs.\ (34)--(37)$\}\cup\{I\}$}{run Algorithm F2$(\psi,g_{\mathrm{block}})$\;}
\end{algorithm}

%==============================================================================
\section{Inverse problem: automorphism design}
\label{sec:inverse}
%==============================================================================

\pq{\S3.4 (section header only; body not yet drafted in the outline)}{%
``3.4 Automorphism Design'' [the outline names this problem but the section
itself is currently empty in the draft --- Algorithm I1 below is this
document's proposed constructive solution, justified by Cor.\ 3.6 and Thm.\
3.5, both already quoted.]
}

\subsection{Algorithm I1 --- orbit-closure construction}

\pq{Thm.\ 3.5 (restated in \S4, ``Notation and Theorems to remember'')}{%
``Theorem: Matching. [\ldots] $\sigma(M)=M \iff \sigma(g)=w\cdot g$ for some
$w\in R^\times$.''
}

\pq{Cor.\ 3.6 (quoted above)}{%
``$\phi$ is a code automorphism of the code if $\phi((f_1,f_2))=u(f_1,f_2)$,
where $u$ is a unit in $R_{\ell,m}^\times$.''
}

\begin{mynote}
Cor.\ 3.6 asks for \emph{some} unit $u$. Algorithm I1 satisfies it with the
simplest possible unit, $u=1$, by construction: if $\phi$ literally fixes
$A$ as a set of monomials ($\phi(A)=A$), the matching equation holds with
$u=1$ trivially. Building such an $A$ requires no search: take $A$ to be a
union of full orbits of the target automorphism group $\Gamma$ (the group
generated by the catalog automorphisms one wants as symmetries) acting on
$\mathbb{Z}_\ell\times\mathbb{Z}_m$.
\end{mynote}

\pq{Prop.\ 2.7 (``Star pairing'')}{%
``The antipode permutes $C$ by the involution $c\mapsto c^*$, $(\alpha,
\beta)\mapsto(\alpha^{-1},\beta^{-1})$, well defined on Frobenius
orbits.''
}

\begin{mynote}
\textbf{Not a direct paper quote.} Prop.\ 2.7 describes how the antipode
acts on the CRT component set $C$; it does not, in the outline as drafted,
state the fact Algorithm I1 actually needs, which is a fact about the
\emph{catalog automorphisms themselves}: every catalog map in \S2.2 (items
1--5) is $\mathbb{Z}$-linear on the exponent lattice $\mathbb{Z}_\ell\times
\mathbb{Z}_m$, and $\overleftarrow{(\cdot)}=-\mathrm{Id}$ is central in
$GL_2(\mathbb{Z})$, so $\phi\circ\overleftarrow{(\cdot)} =
\overleftarrow{(\cdot)}\circ\phi$ for every catalog $\phi$. This is the
author's own supporting lemma, consistent in spirit with Prop.\ 2.7 but
distinct from it, and it is what licenses setting $B:=\overleftarrow{A}$ to
get a second (Hadamard-fold) symmetry for free alongside whatever $\Gamma$
was used to build $A$.
\end{mynote}

\begin{algorithm}[H]
\caption{Orbit-closure construction (constructive Cor.\ 3.6, $u=1$)}
\SetKwInOut{Input}{Input}\SetKwInOut{Output}{Output}
\Input{target generators $\varphi_1,\ldots,\varphi_k$; seed monomials}
\Output{code $(A,B)$ exactly invariant under $\Gamma=\langle\varphi_1,\ldots,\varphi_k\rangle$}
$\Gamma \gets$ BFS closure of $\{\varphi_1,\ldots,\varphi_k\}$ \tcp*{identical routine to Algorithm F3's catalog closure}
\ForEach{seed $(i,j)$}{compute orbit $\{\gamma(i,j):\gamma\in\Gamma\}$\;}
$A \gets$ sum (in $\mathbb F_2$) of the orbit supports\;
\tcp{by construction, $\gamma(A)=A$ for every $\gamma\in\Gamma$ [Cor.\ 3.6, $u=1$]}
\uIf{Hadamard-fold symmetry also wanted}{$B \gets \overleftarrow{A}$ \tcp*{free, by the lemma above}}
\Else{repeat orbit construction with a fresh seed for $B$\;}
\Return{$(A,B)$}
\end{algorithm}

\subsection{Algorithm I2 --- independent re-verification}

\begin{algorithm}[H]
\caption{Re-verification of the constructed code}
\SetKwInOut{Input}{Input}\SetKwInOut{Output}{Output}
\Input{constructed $(A,B)$; group $\Gamma$ from Algorithm I1}
\Output{count of $\Gamma$-elements confirmed as automorphisms}
$n\gets 0$\;
\ForEach{$\gamma\in\Gamma$}{
  \If{Algorithm F1$(\gamma,\mathrm{swap}{=}0)$ returns \textsc{true}}{$n\gets n+1$\;}
}
\Return{$n$ (assert $n=|\Gamma|$)}
\end{algorithm}

\begin{mynote}
Algorithm I2 deliberately reuses Algorithm F1 verbatim rather than trusting
Algorithm I1's construction. This closes the loop without being circular:
Algorithm F1 has no knowledge of how $(A,B)$ was built, so a pass is
independent evidence, not restated assumption.
\end{mynote}

%==============================================================================
\section{Where shared algebraic structure is used}
\label{sec:shared}
%==============================================================================

\begin{enumerate}[leftmargin=*]

\item \textbf{Corollary 2.5's ``(H0)'' licenses both Thm.\ 3.1 and Thm.\
3.2 simultaneously.} Both proofs quoted above invoke Nakayama
\emph{componentwise, because $J_c$ is local even when $R_{\ell,m}$ is not}
--- precisely Cor.\ 2.5. Algorithms F1 and F2 never construct the
components $J_c$ explicitly, but their validity rests on the same single
structural fact: it is what makes ``row space equal $\Rightarrow$ generators
differ by a global unit'' true in the first place, for \emph{both} the
$n=1$ and $n=4$ generator cases.

\item \textbf{One routine, \texttt{generate\_group}, is used by Algorithm
F3 and Algorithm I1.} This was not always true: an earlier version of this
toolkit gave the forward-scan catalog its own fixed composition template
(``one multiplier, one fold, one shear''), which is \emph{not} closed under
composition and under-covers the true group generated by the \S2.2 catalog
items whenever $\ell=m$ (missing interleaved $\tau$-shear compositions) or
whenever a shear is present (the ``(and transposes)'' clause of item 4 was
never wired in at all). Unifying the two call sites onto one BFS-closure
routine simultaneously fixed the forward scan's coverage and confirmed the
inverse-design code's group machinery was correct all along.

\item \textbf{$\mathrm{rank}(H_X)$ and $S=\mathrm{blockdiag}(H_X,H_Z)$ are
computed once per code}, not once per catalog element, and reused as the
fixed operand across the whole scan (Algorithm F3's inner loops).

\item \textbf{The circulant blocks $\rho(A),\rho(B),\rho(\overleftarrow{B}),
\rho(\overleftarrow{A})$} are computed once (\texttt{BBCode.\_\_post\_init\_\_})
and shared across $H_X$, $H_Z$, and $S$ --- the same generator data
underlies every detection level.

\item \textbf{The $\psi_L=\psi_R$ restriction is shared, not just stated.}
\pq{\S3}{``Establish only considering $\psi_L=\psi_R$ in this work\ldots''}
Algorithm F2 applies \emph{one} permutation to all four column-blocks of
$S$ rather than two independent ones, which is both this stated scope
restriction and a direct 4$\times$ reduction in per-test cost.

\end{enumerate}

%==============================================================================
\section{Honest scope}
%==============================================================================

\begin{itemize}[leftmargin=*]
\item The CRT component decomposition itself ($J_c$, Hensel lifts, local
coordinates of \S3.2.1) is never constructed; Lemma~\ref{lem:rowspace}
answers the same yes/no question without it, at the cost of never producing
the per-component unit $\omega_c$ or transport map $T^*_\psi$ the outline's
own proof machinery is built on.
\item $\psi_L\ne\psi_R$ (\S3.3) is out of scope, matching the outline's own
stated restriction.
\item $g_{\mathrm{block}}$ is restricted to the six scalar matrices of
eqs.\ (34)--(37); the fully general $g_{\mathrm{block}}\in M_4(R_{\ell,m})$
of Thm.\ 3.2 is not searched.
\end{itemize}

\end{document}
